% Generated draft: requires mathematical and source review.
\documentclass[11pt]{article}
\usepackage[margin=25mm]{geometry}
\usepackage{fontspec}
\setmainfont{Latin Modern Roman}
\usepackage{amsmath,amssymb,mathtools,unicode-math}
\setmathfont{Latin Modern Math}
\usepackage{enumitem,needspace,xurl,hyperref}
\setlength{\parindent}{0pt}
\setlength{\parskip}{5pt}
\setlength{\emergencystretch}{3em}
\newcommand{\N}{\mathbb{N}}
\newcommand{\Z}{\mathbb{Z}}
\newcommand{\Q}{\mathbb{Q}}
\newcommand{\R}{\mathbb{R}}
\newcommand{\C}{\mathbb{C}}
\newcommand{\F}{\mathbb{F}}
\newcommand{\A}{\mathbb{A}}
\newcommand{\PP}{\mathbb{P}}
\newcommand{\E}{\mathbb{E}}
\newcommand{\eps}{\varepsilon}
\newcommand{\dd}{\,\mathrm{d}}
\begin{document}
\section*{UnsolvedMath ID: 20002814}
\textbf{Source problem code:} catalogue entry
\subsection{Definitions and mathematical statement}
\textbf{Problem 1.2.} Construct a nested supercharacter theory realizing the graded Hopf algebra \(\operatorname{NCQSym}\), indexed by ordered set partitions. The source explicitly says ``nested'' means a projective system of groups, and refers to H. Li's tower-of-algebras work. A complete realization must specify finite groups, supercharacter theories, compatible maps and the product/coproduct operations; a projective system alone does not define those operations.

A supercharacter theory gives a space of superclass functions, with superclass-indicator basis. In the aid's restricted proposed realization, let this basis in degree \(n\) be \(K_\Phi\), indexed by ordered set partitions \(\Phi=(B_1,\ldots,B_k)\) of \([n]\), and require \(K_\Phi\mapsto M_\Phi\), the monomial basis of \(\operatorname{NCQSym}\). The blocks are nonempty, disjoint and ordered. One monomial convention is to sum words whose equal-letter positions are the blocks and whose distinct letter indices increase in the displayed block order.

\textbf{Reported exact coproduct constraint.} For a sublist of blocks, \(\operatorname{std}\) replaces its underlying ordered subset of \([n]\) by consecutive labels, preserving block order. Set the empty list to the degree-zero unit. The required monomial coproduct is
\[
 \Delta M_\Phi=\sum_{r=0}^k
 M_{\operatorname{std}(B_1|\cdots|B_r)}
 \otimes
 M_{\operatorname{std}(B_{r+1}|\cdots|B_k)}.
\]
Consequently, under the prescribed basis map, coalgebra compatibility requires exactly this formula for \(\Delta K_\Phi\), with coefficients one. Its \((a,n-a)\) component is nonzero precisely when
\[
 a=|B_1|+\cdots+|B_r|
\]
for a block prefix. Support alone is only a necessary condition unless coefficients and standardizations also agree; a full Hopf isomorphism additionally needs product/unit/counit compatibility.

For \(\Phi=13|2=(\{1,3\},\{2\})\), the internal term is
\[
 M_{12}\otimes M_1
\]
in bidegree \((2,1)\), while bidegree \((1,2)\) is zero. A candidate restriction coproduct symmetric in all unions of blocks, with nonzero all-subset terms, would also yield \(K_1\otimes K_{12}\), from selecting the second block first. This forbidden term prevents the prescribed basis map from being a coalgebra isomorphism. Thus a direct block-subset-symmetric refinement of the known \(\operatorname{NCSym}\) construction cannot realize this \(\operatorname{NCQSym}\) monomial coproduct unchanged.

This is a necessary design obstruction for the specified basis and restriction template, not a nonexistence theorem for every nested supercharacter theory, a different basis change or modified operations. No actual group/algebra categorification is constructed in the supplied aid. It reports that a 2012 BIRS report explicitly left the problem open and that no later solution was found in its search; the full report reference and search date/scope are not supplied. Those are qualified source reports rather than an independent current-status claim.
\subsection{Short English statement}
A basis-preserving NCQSym realization must split ordered blocks only at prefixes. The degree-three partition \(13|2\) has a \((2,1)\) coproduct term but no \((1,2)\) term, excluding an unchanged restriction rule symmetric over all block subsets. This restricts one template and does not construct or rule out every nested realization.
\subsection{Sources}
\begin{itemize}
\item Local UnsolvedMath registry, permanent ID \(20002814\), catalogue entry, \texttt{aim-representation-theory-notes.json}, \texttt{notes[85]}, \emph{Supercharacters and combinatorial Hopf algebras}, May 17--21, 2010, Problem 1.2: \url{https://aimath.org/WWN/supercharacters/supercharacters.pdf}.
\item Original remark: H. Li, \emph{Algebraic structures on Grothendieck groups of a tower of algebras}; no publication details are supplied. The remark defines nested as a projective system of groups.
\item Supplied research summary invokes an unspecified 2012 BIRS report and gives the prefix-coproduct obstruction. Attempt 1: partial\_result, status at run ``open'', rigor ``proved'', low novelty confidence; registry status partially\_solved. The machine-generated aid is not peer reviewed and requires independent verification before citation; no PDF visit or updated literature search is claimed.
\end{itemize}
\end{document}
